\chapter{Abstract}

\noindent
\textbf{Name:} Max Klemens\\
\textbf{Status:} PhD Inception, first year\\
\textbf{Supervisor:} Prof. Dr. Christian Seidel\\

\vspace{0.5cm}

\section*{Multimodal Airborne Perception and Detect-and-Avoid for Non-Cooperative Drones}

The increasing operation of unmanned aerial vehicles in low-altitude airspace introduces additional risks for crewed aviation, particularly when small drones operate without cooperative identification or communication systems. Conventional detect-and-avoid approaches rely on ground-based sensors or cooperative information sources such as transponders. These approaches are insufficient for the reliable detection of small, non-cooperative drones from a moving airborne platform. The objective of this doctoral research is therefore to investigate multimodal airborne perception and detect-and-avoid methods for helicopter encounters with non-cooperative unmanned aerial vehicles.

The research is structured into three consecutive scientific contributions. The first contribution addresses the development of a synchronized air-to-air RGB--LiDAR dataset and benchmarking framework for the detection, localization, and tracking of small aerial targets. The dataset is intended to include real and simulated sensor data. Particular emphasis will be placed on varying target count (up to six target drones), backgrounds and environmental conditions.

The second contribution will investigate data-driven methods for LiDAR-based detection and temporal tracking of small aerial objects. Existing three-dimensional detection methods are primarily designed for road-based environments and commonly assume objects located close to a ground plane. These assumptions are not directly transferable to unrestricted airspace, where targets may occur at arbitrary elevations and may be represented by only a small number of LiDAR measurements. The planned work will therefore examine suitable spatial representations, temporal point-cloud integration, motion information, and tracking strategies for sparse aerial targets.

The third contribution will integrate the developed perception and tracking methods into a helicopter-oriented detect-and-avoid framework. Representative encounter scenarios will be evaluated under different trajectories, visibility conditions, weather conditions, and sensor degradations. 

The research is currently in the requirements-definition phase. Initial work includes the integration of LiDAR sensing, optical reference tracking, simulation, data recording, and annotation components.

Feedback is specifically requested regarding the scientific delimitation of the three contributions, and the novelty and scope of the proposed airborne benchmark.

